Under the stated allocation,
\[
 \begin{aligned}
 \mathbb E(Z_i\mid t_i)
 &=\frac12(2t_{i,1}-1)-\frac14(2t_{i,2}-1)-\frac14(2t_{i,3}-1)\\
 &=t_{i,1}-\frac{t_{i,2}+t_{i,3}}2=\mu_i.
 \end{aligned}
\]
Because \(Z_i^2=1\), its conditional variance is \(1-\mu_i^2\).  The assignments are independent across units, and \(n^{-1}\sum_i\mu_i=\tau_{c^\dagger}(z)\), so summing the variances proves the mean-squared-error identity.

We next solve the three-observation scalar game exactly.  Let \(N_+=\sum_{i=1}^3\mathbf1\{Z_i=1\}\), \(M_\mu=\sum_i\mu_i\), and \(Q_\mu=\sum_i\mu_i^2\).  Put
\[
 \alpha=\frac{3-\sqrt3}{6},
 \qquad \delta_Z=\alpha(2N_+-3).
\]
Since \(2N_+-3=\sum_iZ_i\), independence and the preceding moment identities give
\[
 \begin{aligned}
 \mathbb E\left\{\delta_Z-\frac{M_\mu}{3}\right\}^2
 &=\alpha^2(3-Q_\mu)+(\alpha-1/3)^2M_\mu^2\\
 &=\frac{2-\sqrt3}{18}\{9-3Q_\mu+M_\mu^2\}\\
 &\leq\frac{2-\sqrt3}{2}=1-\frac{\sqrt3}{2},
 \end{aligned}
\]
where the last step is Cauchy--Schwarz, \(M_\mu^2\leq3Q_\mu\).

For the reverse inequality, put prior weights
\[
 a=\frac{-11+7\sqrt3}{12},\qquad
 b=\frac{8-4\sqrt3}{3},\qquad
 c=\frac{-5+3\sqrt3}{2}
\]
on the five homogeneous mean vectors \((-1,-1,-1)\), \((-1/2,-1/2,-1/2)\), \((0,0,0)\), \((1/2,1/2,1/2)\), and \((1,1,1)\), with respective weights \(a,b,c,b,a\).  These numbers are positive and sum to one.  Direct binomial expansion shows that the observation probabilities \(D_k\), numerator moments \(N_k\), and posterior means for \(N_+=k\) are
\[
 \begin{array}{c|c|c|c}
 k&D_k&N_k&N_k/D_k\\ \hline
 0&(-1+3\sqrt3)/16&(6-5\sqrt3)/16&-3\alpha\\
 1&(9-3\sqrt3)/16&(-6+3\sqrt3)/16&-\alpha\\
 2&(9-3\sqrt3)/16&(6-3\sqrt3)/16&\alpha\\
 3&(-1+3\sqrt3)/16&(-6+5\sqrt3)/16&3\alpha.
 \end{array}
\]
Thus \(\delta_Z\) is the posterior mean under this prior.  Completing the square gives its Bayes risk
\[
 \sum_jw_j\mu_j^2-\sum_{k=0}^3\frac{N_k^2}{D_k}
 =1-\frac{\sqrt3}{2}.
\]
Maximum risk dominates Bayes risk, so the upper and lower bounds coincide and prove the formula for \(V_3^Z\).

We now give the promised full-data witness.  Define
\[
 r_1=\sum_i\mathbf1\{A_i=1\},\qquad
 s_1=\sum_{i:A_i=1}(2Y_i^{\mathrm{obs}}-1),\qquad
 s_-=\sum_{i:A_i\in\{2,3\}}(1-2Y_i^{\mathrm{obs}}).
\]
Let \(1000\delta_F(r_1,s_1,s_-)\) be the following integer table:
\[
\begin{array}{c|r@{\qquad}c|r}
(0,0,-3)&-623&(0,0,-1)&-183\\
(0,0,1)&183&(0,0,3)&623\\
(1,-1,-2)&-638&(1,-1,0)&-199\\
(1,-1,2)&148&(1,1,-2)&-148\\
(1,1,0)&199&(1,1,2)&638\\
(2,-2,-1)&-648&(2,-2,1)&-210\\
(2,0,-1)&-167&(2,0,1)&167\\
(2,2,-1)&210&(2,2,1)&648\\
(3,-3,0)&-660&(3,-1,0)&-167\\
(3,1,0)&167&(3,3,0)&660.
\end{array}
\]
This table is exhaustive: for \(0\leq r\leq3\), the compatible values are \(s\in\{-r,-r+2,\ldots,r\}\) and \(h\in\{-(3-r),-(3-r)+2,\ldots,3-r\}\), giving twenty triples.

Here is an exact finite certificate for its risk.  For a response type \(t\), define the Laurent polynomial
\[
 F_t(U,V,W)=\frac12UV^{2t_1-1}
 +\frac14W^{1-2t_2}+\frac14W^{1-2t_3}.
\]
For an orbit \(m\in\mathcal M_{3,3}\), the coefficient of \(U^rV^sW^h\) in \(\prod_tF_t(U,V,W)^{m_t}\) is exactly the probability of \((r_1,s_1,s_-)=(r,s,h)\).  Therefore every risk is the explicitly evaluable rational number
\[
 R_m=\sum_{r,s,h}[U^rV^sW^h]\prod_tF_t^{m_t}
 \left\{\delta_F(r,s,h)-\frac13\sum_tm_t
 \left(t_1-\frac{t_2+t_3}{2}\right)\right\}^2.
\]
Expanding the eight factors for the \(\binom{10}{7}=120\) response-count orbits and comparing integer numerators after clearing the common denominator gives
\[
 R_m\leq\frac{511653}{4000000}.
\]
Equality holds exactly at the count vectors \((1,0,0,0,0,0,0,2)\) and \((2,0,0,0,0,0,0,1)\), in lexicographic type order \(000,001,010,011,100,101,110,111\).  As a direct boundary check, for the schedule \((000,111,111)\), whose target is zero, the six nonzero orbit probabilities give
\[
 \begin{aligned}
 R&=\frac18(183/1000)^2+\frac18(638/1000)^2
 +\frac14(199/1000)^2+\frac14(167/1000)^2\\
 &\quad+\frac18(648/1000)^2+\frac18(167/1000)^2
 =\frac{511653}{4000000}.
 \end{aligned}
\]
For a fully rational scalar lower certificate, use weights \((1/10,7/20,1/10,7/20,1/10)\) on the same five homogeneous scalar states.  Their observation masses are \((17,15,15,17)/64\), their numerator moments are \((-219,-63,63,219)/1280\), and their prior second moment is \(3/8\).  Completing the square gives
\[
 \frac38-\sum_{k=0}^3\frac{N_k^2}{D_k}
 =\frac{18213}{136000}>\frac{13}{100}
 >\frac{511653}{4000000}.
\]
The strict gap proves the asserted failure of minimax preservation without floating-point arithmetic.

Finally, each feasible scalar mean belongs to \(\{-1,-1/2,0,1/2,1\}\).  Pointwise,
\[
 \mu^2\geq\frac{|\mu|}{2},
 \qquad
 \mu^2\geq\frac{3|\mu|-1}{2}.
\]
Summing and using \(\sum_i|\mu_i|\geq|M_\mu|\) gives the two lower branches for \(Q_\mu\).  If \(|M_\mu|\leq n/2\), equality is obtained by using \(2|M_\mu|\) entries equal to \(\operatorname{sign}(M_\mu)/2\) and all remaining entries zero.  If \(|M_\mu|\geq n/2\), equality is obtained with \(2|M_\mu|-n\) entries equal to \(\operatorname{sign}(M_\mu)\) and the remaining entries equal to \(\operatorname{sign}(M_\mu)/2\).  This proves the wedge formula, including its face junction at \(|M_\mu|=n/2\).